\clearpage\begingroup\SourceLeaf{401b}\ArithmeticHead{70}{L}{STÄCKEL, GAUSS ALS GEOMETER.}
\renewcommand{\thefootnote}{\arabic{footnote})}\setcounter{footnote}{1}
\begin{center}
\textbf{22.}\par\smallskip
\textls[100]{Das Dreieck.}
\end{center}
\noindent\hspace*{\parindent}Rein geometrisch ist der in die Lehrbücher der Elementargeometrie übergegangene
klassische Beweis für den Satz, dass die drei Höhen des Dreiecks
sich in einem Punkte schneiden (W. IV, S 396); er ist 1810 in den Zusätzen
veröffentlicht worden, die \textsc{Gauss} zu \textsc{Schumachers} Übersetzung der \textit{Géométrie de
position} von \textsc{Carnot} beigesteuert hat (Teil 2, Zusatz II, S. 363) Auf einem
verwandten Gedanken beruht der weniger bekannte Beweis von \textsc{Naudé}; dieser
zeigt, dass das Dreieck der Höhenfusspunkte die Höhen zu Winkelhalbierenden
hat\footnote{\textsc{Ph. Naudé}, \textit{Trigonoscopiae cuiusdam novae conspectus}, Miscellanea Berolinensia, t. V, 1737, S. 10; siehe besonders S. 17.}.
\par
In denselben Zusätzen (Zusatz I, S. 359) hat \textsc{Gauss} mittels der Methoden
der analytischen Geometrie einen merkwürdigen Punkt des Dreiecks
\clearpage\SourceLeaf{402a}\ArithmeticHead{71}{R}{ELEMENTARE UND ANALYTISCHE GEOMETRIE.}\setcounter{footnote}{0}
\noindent nachgewiesen, von dem die Durchschnittspunkte der Höhen, der Mittelsenkrechten
und der Schwerlinien besondere Fälle sind (W. IV, S. 393).
\par
Eine handschriftliche Bemerkung zum Zusatz II lehrt, wie die genannten
Durchschnittspunkte mit den komplexen Zahlen zusammenhängen, die den
Ecken des Dreiecks zugeordnet sind (W. IV, S. 396). Mittels komplexer
Grössen ist sicherlich auch die Lösung der Aufgabe gewonnen worden, die
Lage eines Punktes aus den Verhältnissen seiner Abstände von drei der Lage
nach bekannten Punkten zu finden (W. VIII, S. 303).
\par
Endlich ist noch ein Beweis des Pythagoreischen Lehrsatzes aus dem Jahre
1797 (T. Nr. 81) zu nennen, der auf der Ähnlichkeit von Dreiecken beruht\footnote{Der Beweis von \textsc{Gauss} ist den 96 Beweisen hinzuzufügen, die J. \textsc{Versluys} gesammelt hat: \textit{Zes en negentig bewijzen voor het theorema van Pythagoras}, Amsterdam 1914; von den dort mitgeteilten Beweisen kommt dem \textsc{Gauss}schen am nächsten der von \textsc{Brand}, \textit{Une nouvelle démonstration de Pythagore}, Journal de mathématiques élémentaires, série 5, t. 21, 1897, S. 36.}.
\clearpage\endgroup
